% Part: first-order-logic
% Chapter: sequent-calculus
% Section: soundness-identity

\documentclass[../../../include/open-logic-section]{subfiles}

\begin{document}

\olfileid{fol}{seq}{sid}

\olsection{\usetoken{S}{identity}తో నిర్దుష్టత}

\begin{prop}
తాదాత్మ్యానికి ప్రారంభ సీక్వెంట్లు, నియమాలు గల $\Log{LK}$ నిర్దుష్టమైనది.
\end{prop}

\begin{proof}
${} \Sequent \eq[t][t]$ రూపంలోని ప్రారంభ సీక్వెంట్లు చెల్లుబాటు
అవుతాయి; ఎందుకంటే ప్రతి !!{structure}~$\Struct M$కు
$\Sat{M}{\eq[t][t]}$. ($t$ సంవృత పదమని, అంటే అందులో చరరాశులు లేవని
అనుకుంటున్నాం; అందువల్ల చరరాశి విలువనిర్ణయాలు అప్రస్తుతం.)

ఒక !!a{derivation}లో చివరి నిగమనం $=$ అనుకుందాం. అప్పుడు పూర్వాధారం
$\eq[t_1][t_2], \Gamma \Sequent \Delta, !A(t_1)$, నిష్కర్ష
$\eq[t_1][t_2], \Gamma \Sequent \Delta, !A(t_2)$. ఒక
!!a{structure}~$\Struct M$ను పరిశీలించండి. నిష్కర్ష చెల్లుబాటు అని,
అంటే $\Sat{M}{\eq[t_1][t_2]}$, $\Sat{M}{\Gamma}$ అయితే ఏదో
$!C \in \Delta$కు $\Sat{M}{!C}$ లేదా $\Sat{M}{!A(t_2)}$ అని
చూపాలి.

ఆగమన పరికల్పన ప్రకారం పూర్వాధారం చెల్లుబాటు అవుతుంది. అంటే
$\Sat{M}{\eq[t_1][t_2]}$, $\Sat{M}{\Gamma}$ అయితే, (a) ఏదో
$!C \in \Delta$కు $\Sat{M}{!C}$ లేదా (b) $\Sat{M}{!A(t_1)}$.
(a) సందర్భంలో పని పూర్తయింది. (b) సందర్భాన్ని పరిశీలిద్దాం.
$s(x) = \Value{t_1}{M}$ అయ్యే ఒక చరరాశి విలువనిర్ణయం $s$ను తీసుకోండి.
\olref[syn][ass]{prop:sentence-sat-true} ప్రకారం
$\Sat{M}{!A(t_1)}[s]$. $\varAssign{s}{s}{x}$ కనుక
\olref[syn][ext]{prop:ext-formulas} ప్రకారం $\Sat{M}{!A(x)}[s]$.
$\Sat{M}{\eq[t_1][t_2]}$ కనుక
$\Value{t_1}{M} = \Value{t_2}{M}$; అందువల్ల
$s(x) = \Value{t_2}{M}$. \olref[syn][ext]{prop:ext-formulas}ను మళ్లీ
వర్తింపజేస్తే $\Sat{M}{!A(t_2)}[s]$ కూడా వస్తుంది.
\olref[syn][ass]{prop:sentence-sat-true} ప్రకారం $\Sat{M}{!A(t_2)}$.
\end{proof}

\end{document}
